\subsection{RL Pass Filter}

\subsubsection{Overview}
Constructed from a resistor and inductor in series. It
is a single pole filter. An RL Pass Filter can be broken
down into two types: Low Pass Filter and High Pass.
\begin{enumerate}
  \item \textbf{Low Pass Filter:} Allows low frequency
    signals to pass while attenuating high frequency
    signals.
  \item \textbf{High Pass Filter:} Allows high frequency
    signals to pass while attenuating low frequency
    signals.
\end{enumerate}

\subsubsection{Schematic}
\begin{circuitfig}
  % Paths, nodes and wires:
  \draw (4.25, 6.5) -- (5.25, 6.5);
  \node[shape=rectangle, minimum width=0.715cm, minimum height=0.465cm] (N1) at (5.625, 6.5){} node[anchor=center] at (N1.text){$V_{out}$};
  \draw (4.25, 4.75) to[european resistor, l_={$G_1$}, label distance=0.03cm] (4.25, 6.5);
  \draw (2.75, 8.25) to[sinusoidal voltage source, l={$V_1$}] (2.75, 6.75);
  \draw (2.75, 4.75) to[sinusoidal voltage source, l_={$V_2$}] (2.75, 6.25);
  \node[ground, rotate=-90] at (2.75, 6.75){};
  \node[ground, rotate=-90] at (2.75, 6.25){};
  \draw (2.75, 4.75) -- (4.25, 4.75);
  \draw (2.75, 8.25) -- (4.25, 8.25);
  \node[shape=rectangle, minimum width=2.965cm, minimum height=0.465cm] at (3.5, 4.25){} node[anchor=north, align=center, text width=2.577cm, inner sep=6pt] at (3.5, 4.5){Full Schematic};
  \draw (8.625, 6.5) -- (9.625, 6.5);
  \node[shape=rectangle, minimum width=0.715cm, minimum height=0.465cm] (N2) at (10, 6.5){} node[anchor=center] at (N2.text){$V_{out}$};
  \draw (8.625, 4.75) to[european resistor, l_={$G_1$}, label distance=0.03cm] (8.625, 6.5);
  \draw (7.125, 8.25) to[sinusoidal voltage source, l={$V_1$}] (7.125, 6.75);
  \node[ground, rotate=-90] at (7.125, 6.75){};
  \node[ground, rotate=-90] at (7.125, 4.75){};
  \draw (7.125, 4.75) -- (8.625, 4.75);
  \draw (7.125, 8.25) -- (8.625, 8.25);
  \node[shape=rectangle, minimum width=2.965cm, minimum height=0.465cm] at (7.875, 4.25){} node[anchor=north, align=center, text width=2.577cm, inner sep=6pt] at (7.875, 4.5){Low Pass};
  \draw (13, 6.5) -- (14, 6.5);
  \node[shape=rectangle, minimum width=0.715cm, minimum height=0.465cm] (N3) at (14.375, 6.5){} node[anchor=center] at (N3.text){$V_{out}$};
  \draw (13, 4.75) to[european resistor, l_={$G_1$}, label distance=0.03cm] (13, 6.5);
  \draw (11.5, 4.75) to[sinusoidal voltage source, l_={$V_2$}] (11.5, 6.25);
  \node[ground, rotate=-90] at (11.5, 8.25){};
  \node[ground, rotate=-90] at (11.5, 6.25){};
  \draw (11.5, 4.75) -- (13, 4.75);
  \draw (11.5, 8.25) -- (13, 8.25);
  \node[shape=rectangle, minimum width=2.965cm, minimum height=0.465cm] at (12.25, 4.25){} node[anchor=north, align=center, text width=2.577cm, inner sep=6pt] at (12.25, 4.5){High Pass};
  \draw (4.25, 8.25) to[american inductor, l={$L_1$}] (4.25, 6.5);
  \draw (8.625, 8.25) to[american inductor, l={$L_1$}] (8.625, 6.5);
  \draw (13, 8.25) to[american inductor, l={$L_1$}] (13, 6.5);
\end{circuitfig}

\subsubsection{Transfer Function}
Applying KCL at $V_{out}$ on Full Schematic, we get:
\[
  \frac{V_{out} - V_{1}}{sL_1} + (V_{out} - V_{2}){G_1} = 0
\]
\paragraph{Superposition:}
\begin{enumerate}
  \item Due to \(V_{1}\) alone (\(V_{2} = 0V\)):
    This is the low-pass characteristic of the RL Pass Filter.
    \begin{align*}
      \frac{(V_{out} - V_{1})}{sL_1} + (V_{out} - 0){G_1} &= 0 \\
      \frac{V_{out}}{sL_1} - \frac{V_{1}}{sL_1} + V_{out}G_1 &= 0 \\
      V_{out}(G_1 + \frac{1}{sL_1}) - \frac{V_{1}}{sL_1} &= 0 \\
      V_{out}(G_1 + \frac{1}{sL_1}) &= \frac{V_{1}}{sL_1} \\
      V_{out} &= \frac{\frac{V_{1}}{sL_1}}{(G_1 + \frac{1}{sL_1})} \\
      V_{out} &= \frac{1}{L_1 s G_1 + 1}V_{1} \\
      H_1(s) = \frac{N_1(s)}{D_1(s)} = \frac{V_{out_1}}{V_{1}} &= \frac{1}{L_1 s G_1 + 1}
    \end{align*}
  \item Due to \(V_{2}\) alone (\(V_{1} = 0V\)):
    This is the high-pass characteristic of the RL Pass Filter.
    \begin{align*}
      \frac{(V_{out} - 0)}{sL_1} + (V_{out} - V_{2}){G_1} &= 0 \\
      \frac{V_{out}}{sL_1} + V_{out}G_1 - V_{2}G_1 &= 0 \\
      V_{out}(G_1 + \frac{1}{sL_1}) - V_{2}G_1 &= 0 \\
      V_{out}(G_1 + \frac{1}{sL_1}) &= V_{2}G_1 \\
      V_{out} &= \frac{V_{2}G_1}{(G_1 + \frac{1}{sL_1})} \\
      V_{out} &= \frac{sL_1 G_1}{G_1 s L_1 + 1}V_{2} \\
      H_2(s) = \frac{N_2(s)}{D_2(s)} = \frac{V_{out_2}}{V_{2}} &= \frac{sL_1 G_1}{G_1 s L_1 + 1}
    \end{align*}
\end{enumerate}
\paragraph{Total Response:}
\begin{align*}
  V_{out} &= V_{out_1} + V_{out_2} \\
  V_{out} &= H_1(s) V_{1} + H_2(s) V_{2} \\
  V_{out} = \underbrace{\frac{1}{L_1 s G_1 + 1}V_{1}}_{\text{Low Pass Character}} + \underbrace{\frac{sL_1 G_1}{G_1 s L_1 + 1}V_{2}}_{\text{High Pass Character}}
\end{align*}

\subsubsection{Analysis}
\begin{enumerate}
    \item \textbf{Poles (\(D(s) = 0\)):} \\
      For both low pass \(H_1(s)\) and high pass \(H_2(s)\):
      \begin{align*}
        L_1 s G_1 + 1 &= 0 \\
        L_1 s G_1 &= -1 \\
        s &= -\frac{1}{L_1 G_1}
      \end{align*}
      Negative real pole indicates a stable system. Frequency
      corresponding to this pole:
      \begin{align*}
        s &= -2 \pi f_p \\
        -2 \pi f_p &= -\frac{1}{L_1 G_1} \\
        f_p &= \frac{1}{2 \pi L_1 G_1} = \frac{R_1}{2 \pi L_1} \quad\text{where } R_1 = \frac{1}{G_1}
      \end{align*}
    \item \textbf{Zeros (\(N(s) = 0\)):}
      For low pass \(H_1(s)\):
      \begin{align*}
        N_1(s) &= 0 \\
        1 \neq 0
      \end{align*}
      No zeros for \(H_1(s)\). \\
      For high pass \(H_2(s)\):
      \begin{align*}
        N_2(s) &= 0 \\
        s L_1 G_1 = 0 \\
        s_z &= 0 \quad \text{[let \(s = s_z\)]} \\
        f_z &= \frac{|s_z|}{2\pi} = 0 Hz
      \end{align*}
      Zero at the origin for \(H_2(s)\).
    \item \textbf{DC Gain (\(H(0)\)):} \\
      For low pass \(H_1(s)\):
      \begin{align*}
        H_1(0) &= \frac{1}{L_1 (0) G_1 + 1} = 1 \\
        A_{DC_1} = 1 = 0 dB
      \end{align*}
      Unity DC gain for \(H_1(s)\). \\
      For high pass \(H_2(s)\):
      \begin{align*}
        H_2(0) &= \frac{(0) L_1 G_1}{G_1 (0) L_1 + 1} = 0 \\
        A_{DC_2} = 0 = -\infty dB
      \end{align*}
      Zero DC gain for \(H_2(s)\).
    \item \textbf{Gain-Bandwidth Product (\(GBW\)):}
      For low pass \(H_1(s)\):
      \begin{align*}
        GBW &= |A_{DC} \times f_p| \\
        GBW &= 1 \times \frac{R_1}{2 \pi L_1} \\
            &= \frac{R_1}{2 \pi L_1}
      \end{align*}
      For high pass \(H_2(s)\):
      \begin{align*}
        GBW &= A_{DC} \times f_p \\
        GBW &= 0 \times \frac{R_1}{2 \pi L_1} \\
            &= 0
      \end{align*}
      The high pass filter has no gain at DC, so the GBW which
      is a quantity measured from DC to the cutoff frequency is
      zero.
    \item \textbf{Impedance:}
      Both low pass and high pass filters have the same input
      and output impedance.
      Substituting \(G_1 = \frac{1}{R_1}\) and \(s = j \omega\):
      \paragraph{At Input:}
        \begin{align*}
          Z_{in} &= \frac{1}{G_1} + s L_1 \\
          Z_{in} &= R_1 + j \omega L_1 \\
          Z_{in} = R_1 + j \omega L_1
        \end{align*}
      \paragraph{At Output:}
        \begin{align*}
          Z_{out} &= R_1 || s L_1 \\
          Z_{out} &= \frac{R_1 \cdot j \omega L_1}{R_1 + j \omega L_1} \\
          Z_{out} &= \frac{R_1 j \omega L_1}{R_1 + j \omega L_1} \\
          Z_{out} &= \frac{R_1 j \omega L_1 (R_1 - j \omega L_1)}{(R_1 + j \omega L_1)(R_1 - j \omega L_1)} \\
          Z_{out} &= \frac{j \omega R_1^2 L_1 + \omega^2 R_1 L_1^2}{R_1^2 + (\omega L_1)^2} \\
          Z_{out} = \frac{\omega^2 R_1 L_1^2}{R_1^2 + (\omega L_1)^2} + j \frac{\omega R_1^2 L_1}{R_1^2 + (\omega L_1)^2}
        \end{align*}
\end{enumerate}

\subsubsection{Question}
Design a RL Low Pass Filter with a pole frequency of
\(f_p = 1 kHz\) and a resistor value of \(R_1 = 1 k\Omega\).
\paragraph{Solution:}
\begin{enumerate}
  \item Calculating the inductor value \(L_1\):
    \begin{align*}
      f_p &= \frac{R_1}{2 \pi L_1} \\
      L_1 &= \frac{R_1}{2 \pi f_p} \\
      L_1 &= \frac{1 k\Omega}{2 \pi (1 kHz)} \\
      L_1 &\approx 0.159 H
    \end{align*}
\end{enumerate}
Design a RL High Pass Filter with a cutoff frequency of
\(f_p = 45kHz\) using a resistor of value \(R_1 = 330 \Omega\).
Also compute the input and output impedance at \(f = 10kHz\).
\paragraph{Solution:}
\begin{enumerate}
  \item Calculating the inductor value \(L_1\):
    \begin{align*}
      f_p &= \frac{R_1}{2 \pi L_1} \\
      L_1 &= \frac{R_1}{2 \pi f_p} \\
      L_1 &= \frac{330}{2 \pi \cdot 45000} \\
      L_1 = 1.167 mH
    \end{align*}
  \item Calculate the input and output impedance at
    \(f = 10kHz\):
    \paragraph{At Input:}
      \begin{align*}
        Z_{in} &= R_1 + j \omega L_1 \\
        Z_{in} &= 330 + j 2 \pi (10000) (1.167 \times 10^{-3}) \\
        Z_{in} &= 330 + j 73.3 \\
        Z_{in} = 330 + j 73.3 \Omega
      \end{align*}
    \paragraph{At Output:}
      \begin{align*}
        Z_{out} &= \frac{\omega^2 R_1 L_1^2}{R_1^2 + (\omega L_1)^2} + j \frac{\omega R_1^2 L_1}{R_1^2 + (\omega L_1)^2} \\
        Z_{out} &= \frac{(2 \pi (10000))^2 (330) (1.167 \times 10^{-3})^2}{(330)^2 + (2 \pi (10000) (1.167 \times 10^{-3}))^2} + j \frac{(2 \pi (10000)) (330)^2 (1.167 \times 10^{-3})}{(330)^2 + (2 \pi (10000) (1.167 \times 10^{-3}))^2} \\
        Z_{out} &= 6.94 + j 69.4 \\
        Z_{out} = 6.94 + j 69.4 \Omega
      \end{align*}
\end{enumerate}

